\newlength{\anchoflexlazo}
\setlength{\anchoflexlazo}{\dimexpr\textwidth-2.4cm-3.1cm-2.2cm-8\tabcolsep\relax}

\begin{footnotesize}
\begin{longtable}{L{2.4cm}L{3.1cm}L{2.2cm}L{\anchoflexlazo}}
\caption{Definición preliminar del lazo de control.}\label{tab:lazo_control}\\
\toprule
\textbf{Elemento}&\textbf{Definición}&\textbf{Unidad/rango}&\textbf{Observaciones}\\
\midrule
\endfirsthead
\multicolumn{4}{c}{\tablename\ \thetable{} -- continuación}\\
\toprule
\textbf{Elemento}&\textbf{Definición}&\textbf{Unidad/rango}&\textbf{Observaciones}\\
\midrule
\endhead
\midrule
\multicolumn{4}{r}{\textit{Continúa en la página siguiente}}\\
\endfoot
\bottomrule
\endlastfoot
Variable controlada&Temperatura del agua o solución (PT100 conectado al PLC)&$^{\circ}$C&Desarrollo completo en la Sección~\ref{sec:pv}.\\
Variable manipulada&Potencia entregada al calentador (resistencia 2000\,W vía SSR+PWM)&0--2000\,W / 0--100\,\% ciclo de trabajo&Desarrollo completo en la Sección~\ref{sec:mv}.\\
Perturbaciones&Pérdidas térmicas hacia el ambiente (agrupadas en $UA$), $T_{amb}$, adición de fluido frío entre fases, pérdidas en mangueras externas&$UA=13{,}63$\,W/K (experimental); $T_{amb}=18$--$25\,^{\circ}$C&$UA$ se identificó en el ensayo experimental de la Sección~\ref{sec:montaje_control}; las pérdidas reales son cerca de 3,4 veces las supuestas en el modelo preliminar ($\approx 4$\,W/K). La adición de fluido frío es la perturbación más crítica y descarta el control puramente P.\\
Limitaciones&Saturación por sobredimensionamiento; actuador unidireccional&---&Obliga a usar anti-windup en PID\_Compact; el enfriamiento solo ocurre por pérdidas naturales, mucho más lento que el calentamiento.\\
\end{longtable}
\end{footnotesize}